% Generated from THESIS_DRAFT.md.
\chapter{Mở đầu}

\section{Bối cảnh nghiên cứu}

Mô hình thị giác–ngôn ngữ–hành động, gọi tắt là VLA, kết nối quan sát hình ảnh, chỉ dẫn ngôn ngữ và hành động điều khiển robot trong một policy. Khả năng tái sử dụng biểu diễn đã học từ dữ liệu thị giác và ngôn ngữ làm cho hướng tiếp cận này phù hợp với các tác vụ thao tác được mô tả bằng lời. SmolVLA là một đại diện hướng đến mô hình có chi phí triển khai thấp hơn, sử dụng một backbone thị giác–ngôn ngữ cùng action expert để sinh chuỗi hành động \cite{smolvla2025}. Tuy nhiên, một policy có thể thực hiện thao tác đơn lẻ vẫn chưa chắc xử lý được nhiệm vụ mà quyết định tiếp theo phụ thuộc vào lịch sử.

Xét yêu cầu nhấc một chiếc bát khỏi đĩa rồi đặt lại đúng ba lần. Sau mỗi lần đặt lại, hình ảnh bát trên đĩa có thể gần giống ban đầu. Nếu chỉ quan sát cảnh hiện tại, hệ thống khó phân biệt đã hoàn thành một, hai hay ba lần. Hành động cần thực hiện vì thế không chỉ phụ thuộc vị trí hiện tại của bát mà còn phụ thuộc những lần tương tác trước đó. Ví dụ này cũng chỉ ra rằng chiều dài trajectory và mức cần thiết của memory là hai thuộc tính khác nhau: một nhiệm vụ dài có thể bộc lộ tiến độ ngay trong cảnh, trong khi một nhiệm vụ lặp ngắn hơn có thể cần lịch sử để phân biệt các tình huống.

Các nghiên cứu gần đây đã khảo sát nhiều biểu diễn memory. MEM phối hợp lịch sử ảnh ngắn với tóm tắt ngôn ngữ dài, còn MemoryVLA lưu và truy xuất các biểu diễn perceptual và cognitive trong một memory bank \cite{mem2026,memoryvla2026}. LIBERO-Mem xây dựng các tình huống thao tác có thông tin quan trọng gắn với lịch sử của thực thể \cite{liberomem2025}. Những công trình này tạo động lực nghiên cứu, đồng thời cho thấy không thể đánh giá memory chỉ bằng việc thêm một chuỗi văn bản vào instruction rồi quan sát một vài rollout.

\section{Bài toán đặt ra}

Khóa luận nghiên cứu cách bổ sung một bộ nhớ sự kiện bên ngoài SmolVLA. Bộ nhớ phải được xây dựng từ quan sát được phép sử dụng, giữ thông tin có ích qua thời gian và tạo đầu vào phù hợp cho policy. Hai yêu cầu cần được giải quyết đồng thời là tính đúng đắn của thông tin được ghi và khả năng policy khai thác thông tin đó. Một hệ thống ghi nhiều sự kiện nhưng thường nhầm lẫn không đương nhiên tốt hơn hệ thống có khả năng từ chối xác nhận khi thiếu bằng chứng.

Trong phạm vi hiện tại, bài toán được thu hẹp vào quan hệ giữa một chiếc bát và đĩa gốc của nó. Hệ thống theo dõi các lần bát rời đĩa và quay lại. T1 yêu cầu một lần nhấc–đặt lại; T3 yêu cầu ba lần. Đây là tập con của hướng nghiên cứu về thao tác có lịch sử thực thể, chưa đại diện cho memory nhiều vật thể, quan hệ chứa, thứ tự tương tác phức tạp hoặc chuyển giao sang robot thật.

Một khó khăn độc lập với thiết kế memory là tính hợp lệ của phép đo. Policy, simulator, dữ liệu demonstration, bộ tiền xử lý và bộ chấm điểm phải có nguồn gốc rõ ràng. Nếu quy ước ảnh hoặc điều kiện kết thúc khác nhau, chênh lệch success có thể phản ánh thay đổi của đường thực nghiệm thay vì thông tin lịch sử. Vì vậy, xây dựng một phép so sánh có thể kiểm tra lại là một phần của phương pháp nghiên cứu.

Có thể mô tả bài toán ở mức trừu tượng như sau. Tại thời điểm $t$, môi trường có trạng thái đầy đủ $x_t$, nhưng policy chỉ nhận quan sát được cho phép $o_t$, instruction $ell$ và snapshot lịch sử $m_t$. Mục tiêu của task là số chu kỳ cần hoàn thành $N\in\{1,3\}$. Gọi $z_t$ là số chu kỳ vật lý đã thực hiện và $\hat z_t$ là số chu kỳ writer đã xác nhận. Hai đại lượng này không đồng nhất: $z_t$ thuộc trạng thái mô phỏng và evaluator, còn $\hat z_t$ là ước lượng từ bằng chứng RGB có thể thiếu hoặc sai.

Tại mỗi lần truy vấn, policy nhận chuỗi điều kiện $c_t=R(\ell,m_t)$ và sinh một action chunk $A_t$. Yêu cầu thiết kế là duy trì một bộ nhớ có ích cho quyết định nhưng tuân theo ràng buộc thông tin và giao diện:

\begin{equation}
\begin{aligned}
m_t &= U(m_{t-1},o_t,a_{t-1}),\\
A_t &\sim \pi_{\theta,\phi}(\cdot\mid o_t,c_t),\\
\operatorname{Tok}(\ell\oplus R(\ell,m_t)) &\le 48.
\end{aligned}
\end{equation}

Ở đây $a_{t-1}$ là action đã thực thi trước quan sát hiện tại; $\oplus$ biểu thị phép ghép instruction với memory theo đúng quy ước tiền xử lý; $\operatorname{Tok}$ là phép đếm bằng tokenizer checkpoint. Mọi sự kiện được dùng tại thời điểm $t$ phải có nguồn quan sát hợp lệ và thời điểm sẵn sàng không vượt quá $t$. Việc tăng SR không được giả định từ định nghĩa bài toán; nó là một giả thuyết cần kiểm tra bằng rollout có đối chứng.

\section{Mục tiêu và câu hỏi nghiên cứu}

Mục tiêu tổng quát là khảo sát khả năng hỗ trợ tác vụ chuỗi dài của VLA bằng bộ nhớ sự kiện thị giác. Mục tiêu thực hiện ở giai đoạn này gồm thiết kế bộ nhớ, tích hợp với checkpoint có nguồn gốc xác định, huấn luyện thích ứng nhẹ và thực hiện đánh giá sơ bộ trong mô phỏng.

RQ1 đặt câu hỏi liệu có thể đưa memory từ RGB vào SmolVLA trong khi duy trì ranh giới thông tin và contract của policy. RQ2 xem xét liệu lịch sử giúp xử lý tác vụ lặp tốt hơn khi đã kiểm soát tác động của adaptation. RQ3 kiểm tra khả năng duy trì kỹ năng cơ sở sau thích ứng. RQ4 tìm hiểu những giới hạn do chất lượng writer, quy ước dữ liệu và chi phí thực thi gây ra. Các câu hỏi được tách ra để một bằng chứng về khả năng vận hành không bị diễn giải thành bằng chứng về hiệu quả hay tính tổng quát.

Trong phạm vi bản thảo này, RQ1 được trả lời ở mức tích hợp cho hai task và có các kiểm tra riêng về writer trên demonstration. RQ2–RQ4 chưa có đủ phép đo để kết luận. Đặc biệt, so sánh B0–F1 hiện có thay đổi đồng thời trọng số adapter và đường perception–memory; đây là phép đo hệ thống tổng hợp, không nhận diện tác động riêng của lịch sử.

\section{Phạm vi và đóng góp}

Nghiên cứu sử dụng nguồn task và simulator từ LIBERO cùng nhánh LIBERO-Mem đã được khóa revision \cite{liberoSource,liberomemSource}. Giai đoạn V1 thiết lập baseline policy không memory. Phần memory thuộc V2; các kết quả hai giai đoạn không được gộp khi runtime hoặc cadence thực thi khác nhau. Mọi kết quả local trong bản thảo được phân loại là LIBERO-derived và không so ngang với điểm của paper chính thức.

Đóng góp hiện có thứ nhất là một thiết kế tích hợp bộ nhớ sự kiện có bằng chứng và thời điểm khả dụng, gắn với quy tắc xử lý lịch sử chưa đầy đủ. Thứ hai là quy trình nối bộ nhớ với SmolVLA bằng văn bản và thích ứng action expert, có lưu provenance của dữ liệu, tokenizer và adapter. Thứ ba là tập artifact thực nghiệm và cách phân tích giúp phân biệt tiến độ triển khai với mức độ kiểm chứng giả thuyết. Những thành quả này có ý nghĩa ở phương diện thiết kế hệ thống nghiên cứu và thực hành thực nghiệm.

Ở mức thực nghiệm thành phần, audit writer ngày 26/10 ghi nhận 12 lần return trong sáu demonstration được ghép một-một với milestone annotation trong dung sai ±20 frame. Đây là bằng chứng ban đầu rằng pipeline thị giác có thể tạo chuỗi sự kiện đúng thứ tự trong các trajectory đã khảo sát. Các demonstration đều thành công, chỉ đại diện ba trạng thái khởi tạo, và milestone không phải nhãn contact; do đó kết quả này chưa ước lượng được lỗi trên trajectory thất bại hay độ đúng trên phân phối rộng.

Khóa luận kế thừa SmolVLA, Grounding DINO, LoRA, loss huấn luyện và các tài nguyên benchmark. Việc kết hợp chúng không tự chứng minh một nguyên lý memory mới hoặc hiệu năng vượt các công trình liên quan. Đóng góp về cải thiện policy chỉ được xác nhận khi có đối chứng phù hợp. Ở tiến độ hiện tại, kết luận đó vẫn còn mở.

\section{Bố cục khóa luận}

Chương 2 trình bày những khái niệm cần thiết và định vị phương pháp trong các nghiên cứu về memory cho robot. Chương 3 mô tả bộ nhớ sự kiện, giao diện policy, dữ liệu thích ứng và protocol đánh giá. Chương 4 trình bày các thí nghiệm có artifact kiểm tra được, phân tích kết quả và chỉ rõ phần chưa đủ bằng chứng. Chương 5 rút ra kết luận ở phạm vi hiện tại và xác định các bước nghiên cứu tiếp theo.

